\documentclass[12pt,a4paper]{article}
\usepackage[utf8]{vietnam}
\usepackage{amsmath,amssymb}
\usepackage{tabularx}
\usepackage{tikz}
\usepackage{xcolor}
\usepackage[top=0.6cm, bottom=0.6cm, left=0.6cm, right=0.6cm]{geometry}
\usepackage[dethi]{ex_test}

\setlength{\parskip}{0pt}
\setlength{\parindent}{0pt}

% Dòng kẻ chấm mờ đều đặn chuẩn viết tay (khoảng cách 20.5pt)
\newcommand{\dongke}[1]{%
  \begingroup%
  \setlength{\parskip}{0pt}%
  \baselineskip=20.5pt%
  \vspace{1.5mm}%
  \foreach \i in {1,...,#1}{%
    \noindent{\color{gray!60}\dotfill}\par%
  }%
  \vspace{1mm}%
  \endgroup%
}

% Định nghĩa ô bài tập tự luận trong bảng
\newcommand{\oBaiTap}[3][6]{%
  \vspace{1.5mm}%
  \noindent\textbf{#2} #3\par%
  \dongke{#1}%
}

% Câu trắc nghiệm 1 cột: 4 đáp án dàn hàng ngang trên 1 dòng
\newcommand{\cauTN}[6]{% \cauTN{số câu}{đề}{A}{B}{C}{D}
  \noindent\textbf{Câu #1.} #2\par\vspace{0.4mm}%
  \noindent%
  \makebox[0.25\linewidth][l]{\textbf{A.} #3}%
  \makebox[0.25\linewidth][l]{\textbf{B.} #4}%
  \makebox[0.25\linewidth][l]{\textbf{C.} #5}%
  \makebox[0.25\linewidth][l]{\textbf{D.} #6}\par\vspace{0.8mm}%
}

\begin{document}
\pagestyle{empty}

% =============================================================================
% ========================= TOÁN 9 - TUẦN 04 - BUỔI 01: MÃ ĐỀ 101 =============
% ========================= MẶT TRƯỚC (TRANG 1/2) ==============================
% =============================================================================

\noindent
\begin{tabularx}{\linewidth}{|X|c|}
\hline
\textbf{PHIẾU KIỂM TRA ĐỊNH KỲ TOÁN 9 - TUẦN 04 - BUỔI 01} & \textbf{Mã đề: 101 - Trang 1/2} \\
\textit{Chủ đề: Phương trình tích} & \textit{Thời gian: 45 phút} \\
\hline
\end{tabularx}

\vspace{1mm}

\noindent
\begin{tabularx}{\linewidth}{|X|}
\hline
\rule[-3.5mm]{0pt}{10mm}\textbf{Họ và tên học sinh:} \hfill \textbf{Mã HS:} \makebox[3.5cm]{\dotfill} \\
\hline
\end{tabularx}

\vspace{1mm}

\noindent
\begin{tabularx}{\linewidth}{|p{3.2cm}|X|}
\hline
\centering\textbf{ĐIỂM SỐ} & \centering\textbf{LỜI NHẬN XÉT CỦA THẦY CÔ} \tabularnewline
\hline
\rule{0pt}{1.3cm} & \tabularnewline
\hline
\end{tabularx}

\vspace{1.2mm}

\noindent\textbf{BẢNG ĐIỀN ĐÁP ÁN TRẮC NGHIỆM (Mỗi câu đúng 0,5 điểm):}
\vspace{0.8mm}

\noindent
\begingroup
\renewcommand{\arraystretch}{1.35}
\begin{tabularx}{\linewidth}{|c|*{5}{>{\centering\arraybackslash}X|}}
\hline
\textbf{Câu} & \textbf{1} & \textbf{2} & \textbf{3} & \textbf{4} & \textbf{5} \\
\hline
\textbf{Đáp án} & & & & & \\
\hline
\end{tabularx}
\endgroup

\vspace{1.5mm}

\noindent\textbf{\large PHẦN I. TRẮC NGHIỆM KHÁCH QUAN (2,5 điểm)}
\vspace{1mm}

\cauTN{1}{Phương trình nào sau đây là phương trình tích?}{$(2x - 1)(x + 3) = 0$}{$2x - 5 = 3x + 1$}{$x^2 - 4x = 5$}{$\dfrac{2x - 1}{x + 2} = 0$}

\cauTN{2}{Nghiệm của phương trình $(x - 4)(2x + 6) = 0$ là:}{$x = 4$ hoặc $x = 3$}{$x = 4$ hoặc $x = -3$}{$x = -4$ hoặc $x = 3$}{$x = -4$ hoặc $x = -3$}

\cauTN{3}{Tập nghiệm của phương trình $x(x - 7) = 0$ là:}{$S = \{7\}$}{$S = \{0; -7\}$}{$S = \{0; 7\}$}{$S = \{-7\}$}

\cauTN{4}{Số nghiệm của phương trình $(x^2 + 1)(2x - 8) = 0$ là:}{$2$}{$3$}{$1$}{$0$}

\cauTN{5}{Tập nghiệm của phương trình $x^2 - 16 = 0$ là:}{$S = \{4\}$}{$S = \{-4\}$}{$S = \{4; -4\}$}{$S = \{16; -16\}$}

\vspace{1.2mm}

\noindent\textbf{\large PHẦN II. TỰ LUẬN (7,5 điểm)}

\noindent\textbf{Câu 1 (2,5 điểm). Giải các phương trình tích sau:}
\vspace{0.8mm}

\noindent
\begin{tabularx}{\linewidth}{|X|X|}
\hline
\oBaiTap[6]{a)}{$(2x - 5)(3x + 6) = 0$} & 
\oBaiTap[6]{b)}{$(x - 3)(2x + 1) + (x - 3)(x + 5) = 0$} \\
\hline
\oBaiTap[6]{c)}{$(2x + 1)^2 - (x - 2)^2 = 0$} & 
\oBaiTap[6]{d)}{$x^2 - 5x + 6 = 0$} \\
\hline
\end{tabularx}

\newpage
% =============================================================================
% ========================= TOÁN 9 - TUẦN 04 - BUỔI 01: MÃ ĐỀ 101 =============
% ========================= MẶT SAU (TRANG 2/2) ================================
% =============================================================================

\noindent
\begin{tabularx}{\linewidth}{|X|c|}
\hline
\textbf{PHẦN II. TỰ LUẬN (Tiếp theo) - TOÁN 9 TUẦN 04 BUỔI 01} & \textbf{Mã đề: 101 - Trang 2/2} \\
\hline
\end{tabularx}

\vspace{2.5mm}

\noindent\textbf{\large Câu 2 (2,5 điểm). Giải các hệ phương trình sau bằng phương pháp cộng đại số:}
\vspace{0.8mm}

\noindent
\begin{tabularx}{\linewidth}{|X|X|}
\hline
\oBaiTap[12]{a)}{$\begin{cases} 2x + y = 7 \\ 3x - y = 3 \end{cases}$} & 
\oBaiTap[12]{b)}{$\begin{cases} 3x + 2y = 12 \\ 2x - 3y = -5 \end{cases}$} \\
\hline
\end{tabularx}

\vspace{2.5mm}

\noindent\textbf{\large Câu 3 (2,5 điểm). Bài toán thực tế vận chuyển (Trích Đề 10):}
\vspace{0.8mm}

\noindent
\begin{tabularx}{\linewidth}{|X|}
\hline
\vspace{0.5mm}
Để chở $15$ tấn thiết bị phục vụ Lễ kỷ niệm 70 năm chiến thắng Điện Biên Phủ, một đội vận chuyển dự định sử dụng các xe tải loại nhỏ. Do thay đổi kế hoạch, đội quyết định chỉ sử dụng các xe tải loại lớn. Vì vậy, số xe sử dụng giảm đi $2$ xe so với dự định và mỗi xe tải loại lớn chở nhiều hơn mỗi xe tải loại nhỏ là $2$ tấn. Hỏi đội vận chuyển sử dụng bao nhiêu xe tải loại lớn? (Biết mỗi xe tải cùng loại đều chở số tấn thiết bị bằng nhau).\par
\vspace{0.4mm}
\noindent\textit{Bài giải:}\dongke{15}
\\ \hline
\end{tabularx}

\newpage
% =============================================================================
% ========================= TOÁN 9 - TUẦN 04 - BUỔI 01: MÃ ĐỀ 102 =============
% ========================= MẶT TRƯỚC (TRANG 1/2) ==============================
% =============================================================================

\noindent
\begin{tabularx}{\linewidth}{|X|c|}
\hline
\textbf{PHIẾU KIỂM TRA ĐỊNH KỲ TOÁN 9 - TUẦN 04 - BUỔI 01} & \textbf{Mã đề: 102 - Trang 1/2} \\
\textit{Chủ đề: Phương trình tích} & \textit{Thời gian: 45 phút} \\
\hline
\end{tabularx}

\vspace{1mm}

\noindent
\begin{tabularx}{\linewidth}{|X|}
\hline
\rule[-3.5mm]{0pt}{10mm}\textbf{Họ và tên học sinh:} \hfill \textbf{Mã HS:} \makebox[3.5cm]{\dotfill} \\
\hline
\end{tabularx}

\vspace{1mm}

\noindent
\begin{tabularx}{\linewidth}{|p{3.2cm}|X|}
\hline
\centering\textbf{ĐIỂM SỐ} & \centering\textbf{LỜI NHẬN XÉT CỦA THẦY CÔ} \tabularnewline
\hline
\rule{0pt}{1.3cm} & \tabularnewline
\hline
\end{tabularx}

\vspace{1.2mm}

\noindent\textbf{BẢNG ĐIỀN ĐÁP ÁN TRẮC NGHIỆM (Mỗi câu đúng 0,5 điểm):}
\vspace{0.8mm}

\noindent
\begingroup
\renewcommand{\arraystretch}{1.35}
\begin{tabularx}{\linewidth}{|c|*{5}{>{\centering\arraybackslash}X|}}
\hline
\textbf{Câu} & \textbf{1} & \textbf{2} & \textbf{3} & \textbf{4} & \textbf{5} \\
\hline
\textbf{Đáp án} & & & & & \\
\hline
\end{tabularx}
\endgroup

\vspace{1.5mm}

\noindent\textbf{\large PHẦN I. TRẮC NGHIỆM KHÁCH QUAN (2,5 điểm)}
\vspace{1mm}

\cauTN{1}{Phương trình nào sau đây \textbf{không phải} là phương trình tích?}{$(x - 1)(2x + 3) = 0$}{$x(3x - 5) = 0$}{$(x - 2)(x + 4) = 3$}{$(2x + 1)(x - 7) = 0$}

\cauTN{2}{Nghiệm của phương trình $(2x - 8)(x + 5) = 0$ là:}{$x = 4$ hoặc $x = -5$}{$x = -4$ hoặc $x = 5$}{$x = 4$ hoặc $x = 5$}{$x = -4$ hoặc $x = -5$}

\cauTN{3}{Tập nghiệm của phương trình $x(x + 6) = 0$ là:}{$S = \{0; 6\}$}{$S = \{0; -6\}$}{$S = \{-6\}$}{$S = \{6\}$}

\cauTN{4}{Số nghiệm của phương trình $(x^2 + 9)(3x - 6) = 0$ là:}{$1$}{$2$}{$3$}{$0$}

\cauTN{5}{Tập nghiệm của phương trình $x^2 - 25 = 0$ là:}{$S = \{5\}$}{$S = \{-5\}$}{$S = \{5; -5\}$}{$S = \{25; -25\}$}

\vspace{1.2mm}

\noindent\textbf{\large PHẦN II. TỰ LUẬN (7,5 điểm)}

\noindent\textbf{Câu 1 (2,5 điểm). Giải các phương trình tích sau:}
\vspace{0.8mm}

\noindent
\begin{tabularx}{\linewidth}{|X|X|}
\hline
\oBaiTap[6]{a)}{$(3x - 6)(2x + 7) = 0$} & 
\oBaiTap[6]{b)}{$(x + 2)(3x - 1) + (x + 2)(x - 7) = 0$} \\
\hline
\oBaiTap[6]{c)}{$(3x - 1)^2 - (x + 3)^2 = 0$} & 
\oBaiTap[6]{d)}{$x^2 - 7x + 10 = 0$} \\
\hline
\end{tabularx}

\newpage
% =============================================================================
% ========================= TOÁN 9 - TUẦN 04 - BUỔI 01: MÃ ĐỀ 102 =============
% ========================= MẶT SAU (TRANG 2/2) ================================
% =============================================================================

\noindent
\begin{tabularx}{\linewidth}{|X|c|}
\hline
\textbf{PHẦN II. TỰ LUẬN (Tiếp theo) - TOÁN 9 TUẦN 04 BUỔI 01} & \textbf{Mã đề: 102 - Trang 2/2} \\
\hline
\end{tabularx}

\vspace{2.5mm}

\noindent\textbf{\large Câu 2 (2,5 điểm). Giải các hệ phương trình sau bằng phương pháp cộng đại số:}
\vspace{0.8mm}

\noindent
\begin{tabularx}{\linewidth}{|X|X|}
\hline
\oBaiTap[12]{a)}{$\begin{cases} 3x + y = 10 \\ 2x - y = 5 \end{cases}$} & 
\oBaiTap[12]{b)}{$\begin{cases} 2x + 3y = 13 \\ 3x - 2y = 0 \end{cases}$} \\
\hline
\end{tabularx}

\vspace{2.5mm}

\noindent\textbf{\large Câu 3 (2,5 điểm). Bài toán thực tế vận chuyển (Trích Đề 10):}
\vspace{0.8mm}

\noindent
\begin{tabularx}{\linewidth}{|X|}
\hline
\vspace{0.5mm}
Để chở $15$ tấn thiết bị phục vụ Lễ kỷ niệm 70 năm chiến thắng Điện Biên Phủ, một đội vận chuyển dự định sử dụng các xe tải loại nhỏ. Do thay đổi kế hoạch, đội quyết định chỉ sử dụng các xe tải loại lớn. Vì vậy, số xe sử dụng giảm đi $2$ xe so với dự định và mỗi xe tải loại lớn chở nhiều hơn mỗi xe tải loại nhỏ là $2$ tấn. Hỏi đội vận chuyển sử dụng bao nhiêu xe tải loại lớn? (Biết mỗi xe tải cùng loại đều chở số tấn thiết bị bằng nhau).\par
\vspace{0.4mm}
\noindent\textit{Bài giải:}\dongke{15}
\\ \hline
\end{tabularx}

\newpage
% =============================================================================
% ========================= TOÁN 9 - TUẦN 04 - BUỔI 01: MÃ ĐỀ 103 =============
% ========================= MẶT TRƯỚC (TRANG 1/2) ==============================
% =============================================================================

\noindent
\begin{tabularx}{\linewidth}{|X|c|}
\hline
\textbf{PHIẾU KIỂM TRA ĐỊNH KỲ TOÁN 9 - TUẦN 04 - BUỔI 01} & \textbf{Mã đề: 103 - Trang 1/2} \\
\textit{Chủ đề: Phương trình tích} & \textit{Thời gian: 45 phút} \\
\hline
\end{tabularx}

\vspace{1mm}

\noindent
\begin{tabularx}{\linewidth}{|X|}
\hline
\rule[-3.5mm]{0pt}{10mm}\textbf{Họ và tên học sinh:} \hfill \textbf{Mã HS:} \makebox[3.5cm]{\dotfill} \\
\hline
\end{tabularx}

\vspace{1mm}

\noindent
\begin{tabularx}{\linewidth}{|p{3.2cm}|X|}
\hline
\centering\textbf{ĐIỂM SỐ} & \centering\textbf{LỜI NHẬN XÉT CỦA THẦY CÔ} \tabularnewline
\hline
\rule{0pt}{1.3cm} & \tabularnewline
\hline
\end{tabularx}

\vspace{1.2mm}

\noindent\textbf{BẢNG ĐIỀN ĐÁP ÁN TRẮC NGHIỆM (Mỗi câu đúng 0,5 điểm):}
\vspace{0.8mm}

\noindent
\begingroup
\renewcommand{\arraystretch}{1.35}
\begin{tabularx}{\linewidth}{|c|*{5}{>{\centering\arraybackslash}X|}}
\hline
\textbf{Câu} & \textbf{1} & \textbf{2} & \textbf{3} & \textbf{4} & \textbf{5} \\
\hline
\textbf{Đáp án} & & & & & \\
\hline
\end{tabularx}
\endgroup

\vspace{1.5mm}

\noindent\textbf{\large PHẦN I. TRẮC NGHIỆM KHÁCH QUAN (2,5 điểm)}
\vspace{1mm}

\cauTN{1}{Phương trình nào sau đây là phương trình tích?}{$(3x - 2)(x + 4) = 0$}{$3x^2 - 5x + 2 = 1$}{$\dfrac{x + 1}{x - 3} = 0$}{$2x(x - 1) = 6$}

\cauTN{2}{Nghiệm của phương trình $(x + 2)(3x - 9) = 0$ là:}{$x = 2$ hoặc $x = -3$}{$x = -2$ hoặc $x = -3$}{$x = 2$ hoặc $x = 3$}{$x = -2$ hoặc $x = 3$}

\cauTN{3}{Tập nghiệm của phương trình $2x(x - 4) = 0$ là:}{$S = \{4\}$}{$S = \{0; 4\}$}{$S = \{0; -4\}$}{$S = \{2; 4\}$}

\cauTN{4}{Số nghiệm của phương trình $(x^2 + 4)(x + 1)(x - 3) = 0$ là:}{$2$}{$3$}{$4$}{$1$}

\cauTN{5}{Tập nghiệm của phương trình $4x^2 - 9 = 0$ là:}{$S = \left\{\dfrac{3}{2}\right\}$}{$S = \left\{-\dfrac{3}{2}\right\}$}{$S = \left\{\dfrac{3}{2}; -\dfrac{3}{2}\right\}$}{$S = \left\{\dfrac{9}{4}; -\dfrac{9}{4}\right\}$}

\vspace{1.2mm}

\noindent\textbf{\large PHẦN II. TỰ LUẬN (7,5 điểm)}

\noindent\textbf{Câu 1 (2,5 điểm). Giải các phương trình tích sau:}
\vspace{0.8mm}

\noindent
\begin{tabularx}{\linewidth}{|X|X|}
\hline
\oBaiTap[6]{a)}{$(4x - 8)(2x + 3) = 0$} & 
\oBaiTap[6]{b)}{$(x - 4)(3x + 2) + (x - 4)(x - 6) = 0$} \\
\hline
\oBaiTap[6]{c)}{$(2x - 3)^2 - (x + 1)^2 = 0$} & 
\oBaiTap[6]{d)}{$x^2 - 6x + 8 = 0$} \\
\hline
\end{tabularx}

\newpage
% =============================================================================
% ========================= TOÁN 9 - TUẦN 04 - BUỔI 01: MÃ ĐỀ 103 =============
% ========================= MẶT SAU (TRANG 2/2) ================================
% =============================================================================

\noindent
\begin{tabularx}{\linewidth}{|X|c|}
\hline
\textbf{PHẦN II. TỰ LUẬN (Tiếp theo) - TOÁN 9 TUẦN 04 BUỔI 01} & \textbf{Mã đề: 103 - Trang 2/2} \\
\hline
\end{tabularx}

\vspace{2.5mm}

\noindent\textbf{\large Câu 2 (2,5 điểm). Giải các hệ phương trình sau bằng phương pháp cộng đại số:}
\vspace{0.8mm}

\noindent
\begin{tabularx}{\linewidth}{|X|X|}
\hline
\oBaiTap[12]{a)}{$\begin{cases} x + 2y = 8 \\ 3x - 2y = 8 \end{cases}$} & 
\oBaiTap[12]{b)}{$\begin{cases} 3x + 4y = 18 \\ 2x - 3y = -5 \end{cases}$} \\
\hline
\end{tabularx}

\vspace{2.5mm}

\noindent\textbf{\large Câu 3 (2,5 điểm). Bài toán thực tế vận chuyển (Trích Đề 10):}
\vspace{0.8mm}

\noindent
\begin{tabularx}{\linewidth}{|X|}
\hline
\vspace{0.5mm}
Để chở $15$ tấn thiết bị phục vụ Lễ kỷ niệm 70 năm chiến thắng Điện Biên Phủ, một đội vận chuyển dự định sử dụng các xe tải loại nhỏ. Do thay đổi kế hoạch, đội quyết định chỉ sử dụng các xe tải loại lớn. Vì vậy, số xe sử dụng giảm đi $2$ xe so với dự định và mỗi xe tải loại lớn chở nhiều hơn mỗi xe tải loại nhỏ là $2$ tấn. Hỏi đội vận chuyển sử dụng bao nhiêu xe tải loại lớn? (Biết mỗi xe tải cùng loại đều chở số tấn thiết bị bằng nhau).\par
\vspace{0.4mm}
\noindent\textit{Bài giải:}\dongke{15}
\\ \hline
\end{tabularx}

\newpage
% =============================================================================
% ========================= TOÁN 9 - TUẦN 04 - BUỔI 01: MÃ ĐỀ 104 =============
% ========================= MẶT TRƯỚC (TRANG 1/2) ==============================
% =============================================================================

\noindent
\begin{tabularx}{\linewidth}{|X|c|}
\hline
\textbf{PHIẾU KIỂM TRA ĐỊNH KỲ TOÁN 9 - TUẦN 04 - BUỔI 01} & \textbf{Mã đề: 104 - Trang 1/2} \\
\textit{Chủ đề: Phương trình tích} & \textit{Thời gian: 45 phút} \\
\hline
\end{tabularx}

\vspace{1mm}

\noindent
\begin{tabularx}{\linewidth}{|X|}
\hline
\rule[-3.5mm]{0pt}{10mm}\textbf{Họ và tên học sinh:} \hfill \textbf{Mã HS:} \makebox[3.5cm]{\dotfill} \\
\hline
\end{tabularx}

\vspace{1mm}

\noindent
\begin{tabularx}{\linewidth}{|p{3.2cm}|X|}
\hline
\centering\textbf{ĐIỂM SỐ} & \centering\textbf{LỜI NHẬN XÉT CỦA THẦY CÔ} \tabularnewline
\hline
\rule{0pt}{1.3cm} & \tabularnewline
\hline
\end{tabularx}

\vspace{1.2mm}

\noindent\textbf{BẢNG ĐIỀN ĐÁP ÁN TRẮC NGHIỆM (Mỗi câu đúng 0,5 điểm):}
\vspace{0.8mm}

\noindent
\begingroup
\renewcommand{\arraystretch}{1.35}
\begin{tabularx}{\linewidth}{|c|*{5}{>{\centering\arraybackslash}X|}}
\hline
\textbf{Câu} & \textbf{1} & \textbf{2} & \textbf{3} & \textbf{4} & \textbf{5} \\
\hline
\textbf{Đáp án} & & & & & \\
\hline
\end{tabularx}
\endgroup

\vspace{1.5mm}

\noindent\textbf{\large PHẦN I. TRẮC NGHIỆM KHÁCH QUAN (2,5 điểm)}
\vspace{1mm}

\cauTN{1}{Phương trình nào sau đây là phương trình tích?}{$2x^2 + 5x = 7$}{$(x - 5)(4x + 8) = 0$}{$\dfrac{3x - 1}{x + 1} = 0$}{$(x - 1)^2 = 4$}

\cauTN{2}{Nghiệm của phương trình $(2x - 6)(x + 1) = 0$ là:}{$x = -3$ hoặc $x = 1$}{$x = 3$ hoặc $x = 1$}{$x = -3$ hoặc $x = -1$}{$x = 3$ hoặc $x = -1$}

\cauTN{3}{Tập nghiệm của phương trình $3x(x + 5) = 0$ là:}{$S = \{-5\}$}{$S = \{0; -5\}$}{$S = \{0; 5\}$}{$S = \{3; -5\}$}

\cauTN{4}{Số nghiệm của phương trình $(x^2 + 25)(2x - 10) = 0$ là:}{$1$}{$2$}{$3$}{$0$}

\cauTN{5}{Tập nghiệm của phương trình $x^2 - 36 = 0$ là:}{$S = \{6\}$}{$S = \{-6\}$}{$S = \{6; -6\}$}{$S = \{36; -36\}$}

\vspace{1.2mm}

\noindent\textbf{\large PHẦN II. TỰ LUẬN (7,5 điểm)}

\noindent\textbf{Câu 1 (2,5 điểm). Giải các phương trình tích sau:}
\vspace{0.8mm}

\noindent
\begin{tabularx}{\linewidth}{|X|X|}
\hline
\oBaiTap[6]{a)}{$(5x - 10)(3x + 4) = 0$} & 
\oBaiTap[6]{b)}{$(x + 3)(2x - 5) + (x + 3)(x - 1) = 0$} \\
\hline
\oBaiTap[6]{c)}{$(3x + 2)^2 - (x - 4)^2 = 0$} & 
\oBaiTap[6]{d)}{$x^2 - 8x + 12 = 0$} \\
\hline
\end{tabularx}

\newpage
% =============================================================================
% ========================= TOÁN 9 - TUẦN 04 - BUỔI 01: MÃ ĐỀ 104 =============
% ========================= MẶT SAU (TRANG 2/2) ================================
% =============================================================================

\noindent
\begin{tabularx}{\linewidth}{|X|c|}
\hline
\textbf{PHẦN II. TỰ LUẬN (Tiếp theo) - TOÁN 9 TUẦN 04 BUỔI 01} & \textbf{Mã đề: 104 - Trang 2/2} \\
\hline
\end{tabularx}

\vspace{2.5mm}

\noindent\textbf{\large Câu 2 (2,5 điểm). Giải các hệ phương trình sau bằng phương pháp cộng đại số:}
\vspace{0.8mm}

\noindent
\begin{tabularx}{\linewidth}{|X|X|}
\hline
\oBaiTap[12]{a)}{$\begin{cases} 3x + 2y = 13 \\ 5x - 2y = 11 \end{cases}$} & 
\oBaiTap[12]{b)}{$\begin{cases} 4x + 3y = 17 \\ 3x - 2y = 0 \end{cases}$} \\
\hline
\end{tabularx}

\vspace{2.5mm}

\noindent\textbf{\large Câu 3 (2,5 điểm). Bài toán thực tế vận chuyển (Trích Đề 10):}
\vspace{0.8mm}

\noindent
\begin{tabularx}{\linewidth}{|X|}
\hline
\vspace{0.5mm}
Để chở $15$ tấn thiết bị phục vụ Lễ kỷ niệm 70 năm chiến thắng Điện Biên Phủ, một đội vận chuyển dự định sử dụng các xe tải loại nhỏ. Do thay đổi kế hoạch, đội quyết định chỉ sử dụng các xe tải loại lớn. Vì vậy, số xe sử dụng giảm đi $2$ xe so với dự định và mỗi xe tải loại lớn chở nhiều hơn mỗi xe tải loại nhỏ là $2$ tấn. Hỏi đội vận chuyển sử dụng bao nhiêu xe tải loại lớn? (Biết mỗi xe tải cùng loại đều chở số tấn thiết bị bằng nhau).\par
\vspace{0.4mm}
\noindent\textit{Bài giải:}\dongke{15}
\\ \hline
\end{tabularx}

\end{document}
